% Options for packages loaded elsewhere
\PassOptionsToPackage{unicode}{hyperref}
\PassOptionsToPackage{hyphens}{url}
\documentclass[
  a4paper,
]{article}
\usepackage{xcolor}
\usepackage[margin=25mm]{geometry}
\usepackage{amsmath,amssymb}
\setcounter{secnumdepth}{-\maxdimen} % remove section numbering
\usepackage{iftex}
\ifPDFTeX
  \usepackage[T1]{fontenc}
  \usepackage[utf8]{inputenc}
  \usepackage{textcomp} % provide euro and other symbols
\else % if luatex or xetex
  \usepackage{unicode-math} % this also loads fontspec
  \defaultfontfeatures{Scale=MatchLowercase}
  \defaultfontfeatures[\rmfamily]{Ligatures=TeX,Scale=1}
\fi
\usepackage{lmodern}
\ifPDFTeX\else
  % xetex/luatex font selection
\fi
% Use upquote if available, for straight quotes in verbatim environments
\IfFileExists{upquote.sty}{\usepackage{upquote}}{}
\IfFileExists{microtype.sty}{% use microtype if available
  \usepackage[]{microtype}
  \UseMicrotypeSet[protrusion]{basicmath} % disable protrusion for tt fonts
}{}
\makeatletter
\@ifundefined{KOMAClassName}{% if non-KOMA class
  \IfFileExists{parskip.sty}{%
    \usepackage{parskip}
  }{% else
    \setlength{\parindent}{0pt}
    \setlength{\parskip}{6pt plus 2pt minus 1pt}}
}{% if KOMA class
  \KOMAoptions{parskip=half}}
\makeatother
\usepackage{longtable,booktabs,array}
\newcounter{none} % for unnumbered tables
\usepackage{calc} % for calculating minipage widths
% Correct order of tables after \paragraph or \subparagraph
\usepackage{etoolbox}
\makeatletter
\patchcmd\longtable{\par}{\if@noskipsec\mbox{}\fi\par}{}{}
\makeatother
% Allow footnotes in longtable head/foot
\IfFileExists{footnotehyper.sty}{\usepackage{footnotehyper}}{\usepackage{footnote}}
\makesavenoteenv{longtable}
\usepackage{graphicx}
\makeatletter
\newsavebox\pandoc@box
\newcommand*\pandocbounded[1]{% scales image to fit in text height/width
  \sbox\pandoc@box{#1}%
  \Gscale@div\@tempa{\textheight}{\dimexpr\ht\pandoc@box+\dp\pandoc@box\relax}%
  \Gscale@div\@tempb{\linewidth}{\wd\pandoc@box}%
  \ifdim\@tempb\p@<\@tempa\p@\let\@tempa\@tempb\fi% select the smaller of both
  \ifdim\@tempa\p@<\p@\scalebox{\@tempa}{\usebox\pandoc@box}%
  \else\usebox{\pandoc@box}%
  \fi%
}
% Set default figure placement to htbp
\def\fps@figure{htbp}
\makeatother
\setlength{\emergencystretch}{3em} % prevent overfull lines
\providecommand{\tightlist}{%
  \setlength{\itemsep}{0pt}\setlength{\parskip}{0pt}}
\usepackage{bookmark}
\IfFileExists{xurl.sty}{\usepackage{xurl}}{} % add URL line breaks if available
\urlstyle{same}
\hypersetup{
  hidelinks,
  pdfcreator={LaTeX via pandoc}}

\author{}
\date{}

\begin{document}

\section{Onset and Threefold Classification of Outcomes of Spontaneous
Splitting in N-Body Relational-Wave Closed
Systems}\label{onset-and-threefold-classification-of-outcomes-of-spontaneous-splitting-in-n-body-relational-wave-closed-systems}

\subsection{Unfreezing the generator alone makes the parent state decay
exponentially at a rate independent of the seed amplitude, and the
outcomes divide into expansion, bounded mixing, and recurrence ---
introducing the scale-anonymity axiom, with the stopping of splitting as
a connection
problem}\label{unfreezing-the-generator-alone-makes-the-parent-state-decay-exponentially-at-a-rate-independent-of-the-seed-amplitude-and-the-outcomes-divide-into-expansion-bounded-mixing-and-recurrence-introducing-the-scale-anonymity-axiom-with-the-stopping-of-splitting-as-a-connection-problem}

\textbf{Author:} Noriaki Kihara \textbf{ORCID:} 0009-0004-6753-4020
\textbf{Date:} July 22, 2026 \textbf{Version DOI:}
10.5281/zenodo.21486234 \textbf{Concept DOI:} 10.5281/zenodo.21486233
\textbf{Position:} Dimensional Generation Structure series, Paper 5, v1

\begin{center}\rule{0.5\linewidth}{0.5pt}\end{center}

\subsection{Abstract}\label{abstract}

Theorem 5 of Paper 4 {[}3{]} established that under a fixed generator
neither closure quantities nor energy-like quantities move between
rotation planes --- a no-go for spontaneous splitting of the state
(Consequences 16.5 and 16.6 {[}1{]}). Taking this no-go as the starting
boundary, this paper answers the following question.

\begin{quote}
Can multiple waves split spontaneously under the closure condition? If
they split, can the subsequent outcomes be classified?
\end{quote}

To pose the question precisely, we first separate two layers. The
closure \(\sum_k x_k^2=0\) imposes no constraint on the number of terms
\(\mathcal N\) of a decomposition --- any term \(X_P^2\) can always be
rewritten as \(Y_1^2+Y_2^2\) over the complex numbers (always solvable,
with a continuum of solutions), and no obstacle to
\(\mathcal N\to\mathcal N+1\) exists on the axiom side. The number of
waves is not determined inside the system, and \textbf{the layer of
existence has no ceiling (\(\mathcal N\) does not saturate)}. We
separate the symbols: \(\mathcal N\) is the number of terms of a
decomposition (layer of existence), and \(N\) is the number of bodies,
the model parameter of this paper (a structural specification of the
representation; the flat representation has \(M\) relational waves).
What this paper measures is therefore not ``permission to split'' in the
layer of existence. It measures the dynamics of the \textbf{readout
layer} --- register allocation, active planes, freezing. The structures
that look finite across the series (the \(O(N)\) cap on generator rank
{[}2{]}, the saturation of spatial directions at three {[}2{]}, register
freezing {[}3{]}, and the outcome classification of this paper) are all
saturations of the readout, not finiteness of existence.

Three preparations are made. First, we introduce scale anonymity as
Axiom 0.5. The relations, update rules, and readouts of the system may
reference no absolute scale, and every construction must be covariant
under \(Z\to\lambda Z\) (\(\lambda\in\mathbb C^\ast\)). This is an
unprovable formation rule of the same rank as Axiom 0 (anonymity), and
the zero right-hand side of the quadratic closure of Axiom 1 is the
unique choice compatible with this axiom within the family
\(\sum_n x_n^2=c\). Second, we show that every smooth dynamics
preserving the closure sphere is exhausted by the form \(\dot X=K(X)X\)
(\(K\) skew-symmetric) (Theorem 2). The fixed generator is merely the
special case \(K(X)=K_0\), and the search space for ``the minimal
extension enabling splitting'' is thereby closed. Third, as the minimal
extension intrinsic to the model, the phase-difference sine generator
\(K_{ef}=A_{ef}\sin(\theta_f-\theta_e)\) is reconstructed at every step
from the current phases \(\arg Z(\tau)\) (unfreezing). Being a function
of phase differences only, this generator is zeroth-order homogeneous
and satisfies Axiom 0.5 without adjustment.

Spontaneity is defined operationally: the parent state must be
dynamically unstable, i.e., arbitrarily small perturbations must grow
exponentially at an intrinsic rate independent of the perturbation
amplitude. Two structures were fixed in the construction of the parent
state. Real-coefficient in-plane states have collinear component phases,
on which the sine generator vanishes exactly (Theorem 4), so a
self-consistent parent must be the circularly polarized eigenmode
\(v=(p+iq)/\sqrt2\); and since an eigenmode advances by a common phase
rotation only, it is an exact periodic orbit (Proposition 1; numerical
verification: generator variation \(9.2\times10^{-14}\) over one
period). As a numerical insight, the generator of the self-consistent
parent has \(\operatorname{rank}2\) (one rotation plane), and the
remaining \(M-2\) dimensions form a static kernel.

In the numerical experiments (\(N=4\), \(M=6\), seed amplitudes
\(\delta=10^{-2}\)--\(10^{-5}\), 6 trials each, 4320 steps; the dormant
seed is constructed from a real orthonormal pair in the kernel as
\(g=(u+iw)/\sqrt2\) with \(g^{\mathsf T}g=0\) exact, and the initial
state satisfies zero closure to
\(|Z^{\mathsf T}Z|\le3.2\times10^{-13}\)), all three criteria of
spontaneity held. The initial exponential growth rate of the dormant
fraction \(f(\tau)\) was invariant over three decades of \(\delta\),
with medians \(3.78\times10^{-2}\)--\(3.80\times10^{-2}\)/step;
absolute-threshold crossing times were linear in \(\log(1/\delta)\)
(about 84 steps per decade); and relative-threshold crossing times were
\(\delta\)-independent (about 84 steps) (the growth is not
single-exponential: the initial-fit rate and the effective rate derived
from crossing times differ; Section 8.2). Decisive is the \(\delta=0\)
verification. The applied perturbation is zero, and under exact
arithmetic the orbit would be invariant. Nevertheless, the nonzero
numerical perturbation of floating-point rounding (\(f\sim10^{-33}\))
was amplified by about 20 orders of magnitude, reaching \(f>10^{-10}\)
at \(\tau\approx850\) and saturating --- strong evidence of transverse
instability. Throughout, the real and imaginary parts of the complex
bilinear closure (corresponding to the exact splitting conditions of
Paper 4) were conserved to deviations \(\le1.8\times10^{-14}\). In the
control experiment (same initial conditions, fixed generator),
reference-register variation was \(\le4.5\times10^{-11}\): the numerical
reconfirmation of the no-go and the onset of splitting were established
on the same code. Scale anonymity was verified numerically as well:
under the scale transformation \(\lambda=2^{10}\) all ratio observables
agreed bitwise, and for generic \(\lambda\) and common phase rotations
deviations stayed \(\le1.1\times10^{-13}\). The onset law is also robust
in the number of bodies. Runs with the low-rank vertex-decomposition
engine (an algorithmic use of the decomposition \(K=WJW^{\mathsf T}\) of
Paper 3) at \(N=40,300,1000\) (\(M=780,\,44850,\,499500\)), measuring
the dormant fraction without cancellation via projection onto the
orthogonal complement of the parent plane, directly observed
constant-rate exponential amplification over more than 20 decades from
seed depths \(f(0)\approx\delta^2=10^{-30}\) (Experiment O6, Section
8.8). The rate depends only weakly on \(N\)
(\(3.6\times10^{-2}\to2.1\times10^{-2}\)/step, decreasing monotonically
with \(N\) within the same order).

In the outcome-classification experiments (\(N=3,4,5,6\), two families
of initial conditions --- near-parent and generic zero-closure states
--- 6 trials each, 20000 steps), the outcomes divided into three
classes: expansion (a new rotation plane establishes itself and the
allocation spreads over all relational waves), bounded mixing (the
allocation spreads but no second plane establishes itself), and
recurrence (repeated deep returns toward the parent plane). Class
selection has a sharp \(N\) phase structure. At \(N=3\) expansion is
structurally impossible (a skew-symmetric matrix with \(M=3\) has
\(\operatorname{rank}\le2\)); at \(N=4\) zero-closure states did not
expand either (both families bounded mixing); for \(N\ge5\) every
initial condition examined (2 families × 6 seeds) expanded --- the
boundary lies between \(N=4\) and \(N=5\). Furthermore, a controlled
contrast experiment (Experiment O5: everything identical except the
initial value of the complex bilinear closure) confirmed that the ratio
\(c=\lvert Z^{\mathsf T}Z\rvert/\lVert Z\rVert^2\), exactly conserved by
the dynamics, stratifies the outcome classes at small \(N\) --- \(N=4\):
bounded mixing at \(c=0\) (3/3) but expansion at \(c\approx0.6\) (3/3);
\(N=3\): bounded mixing at \(c=0\) (3/3) but recurrence at
\(c\approx0.55\) (3/3). Closure is not only a conservation law; its
conserved value selects the outcome (Section 8.7). In addition, in the
\(N=4\) parent family the second singular direction \(\sigma_2\) decayed
exponentially without external intervention
(\(6\times10^{-4}\to10^{-5}\) over 20000 steps), an observed route by
which the system spontaneously refreezes into a single-plane
self-consistent configuration. The stopping mechanism of splitting,
together with this observation, is the central subject of the sequel.

\begin{center}\rule{0.5\linewidth}{0.5pt}\end{center}

\subsection{0. Conclusions}\label{conclusions}

\[
\boxed{
\begin{aligned}
&\text{Closure-preserving dynamics is exhausted by }\dot X=K(X)X\ (K\text{ skew-symmetric});\text{ the fixed generator is a special case.}\\
&\text{Unfreezing the generator alone makes the parent state (circularly polarized eigenmode) decay exponentially}\\
&\text{without external forcing, while the real and imaginary parts of the closure are conserved to machine precision}\\
&\text{and the allocation moves spontaneously into dormant directions.}\\
&\text{The outcomes divide into three classes — expansion, bounded mixing, recurrence — selected by }N\\
&\text{(structural room) and the conserved ratio }c=\lvert Z^{\mathsf T}Z\rvert/\lVert Z\rVert^2.
\end{aligned}
}
\]

Spontaneity holds as dynamical instability. The growth rate is invariant
over three decades of the seed amplitude \(\delta\), and even the exact
periodic orbit at \(\delta=0\) collapses from rounding noise alone.
Under a deterministic update rule, which direction the state collapses
into is decided by microscopic differences not contained in the
macroscopic description --- determinism and predictability separate in
this minimal model.

The outcome classes have an \(N\) phase structure. At \(N=3\) expansion
is structurally impossible; at \(N=4\) zero-closure states do not
expand; for \(N\ge5\) every initial condition examined expanded --- the
boundary lies between \(N=4\) and \(N=5\). The controlled contrast
(Experiment O5) confirmed that what acts is the exactly conserved ratio
\(c=\lvert Z^{\mathsf T}Z\rvert/\lVert Z\rVert^2\): at \(c=0\) (the
zero-closure stratum = the stratum designated by Axiom 1) expansion is
suppressed for \(N\le4\), while at \(c\approx0.5\)--\(0.6\), \(N=4\)
expands and \(N=3\) recurs. The value of the conserved quantity
stratifies the outcome. In the \(N=4\) parent family the second singular
direction decays exponentially of its own accord, heading toward
refreezing --- evidence that stopping can occur without an external
freezing operation, and the central object of the sequel (stopping of
splitting).

All statements of this paper belong to the readout layer. The layer of
existence (\(\mathcal N\) itself) is unconstrained by the axioms and
does not saturate --- every structure that looks finite is a saturation
of the readout (Section 1.2).

Scale anonymity (Axiom 0.5) runs through the entire construction. The
unfrozen generator is zeroth-order homogeneous as a function of phase
differences only; all readouts are ratio quantities; all activity
thresholds are relative. Numerically, all ratio observables agreed
bitwise under a \(2^{10}\)-fold scale change.

\begin{center}\rule{0.5\linewidth}{0.5pt}\end{center}

\subsection{1. Research Question}\label{research-question}

\subsubsection{1.1 The no-go as the starting
boundary}\label{the-no-go-as-the-starting-boundary}

Paper 3 {[}2{]} proved that the fixed generator of N-body complete
pairwise relational waves decomposes into a direct sum of rotation
planes and a null space, with the number of rotation modes growing at
most linearly. Paper 4 {[}3{]} established that each plane carries
kinematics isomorphic to the two-body AB system, and its Theorem 5 gave
a strong no-go: under a fixed generator, neither the closure quantities
\(Q_j\) nor the energy-like quantities \(H_j\) move between planes or
with the kernel (Consequences 16.5 and 16.6 {[}1{]}).

Spontaneous splitting of the state --- reallocation between registers,
activation of new rotation modes --- therefore cannot occur in the
fixed-generator dynamics as a matter of principle. To ask about
splitting, a relational update rule beyond the fixed generator is
required. This paper identifies the minimal such extension, determines
numerically whether splitting begins, and classifies the outcomes after
onset.

This paper is the first of a three-part sequence. First question (this
paper): can splitting begin? Second question (sequel): can splitting
stop? Third question (sequel): can the conserved quantities generated by
splitting differentiate?

\subsubsection{1.2 Layer of existence and layer of readout --- what this
paper
measures}\label{layer-of-existence-and-layer-of-readout-what-this-paper-measures}

The closure \(\sum_k x_k^2=0\) does not constrain the number of terms
\(\mathcal N\) of a decomposition. The splitting \(X_P^2=Y_1^2+Y_2^2\)
is always solvable over the complex numbers, with a continuum of
solutions --- the splitting map is not a ``hard-to-construct missing
piece'' but an \textbf{underdetermined freedom}. In the layer of
existence the number of waves is arbitrary (it may be infinite) and does
not saturate.

Distinction of symbols: \(\mathcal N\) (number of terms of a
decomposition, layer of existence) and \(N\) (number of bodies,
structural specification of the model = external input) are different
quantities. The \(N\) used by the experiments of this paper for
parameterization is the latter, and the number of terms of its flat
representation is the number of relational waves \(M=\binom N2\).

By contrast, every finiteness the series has established belongs to the
readout layer.

{\def\LTcaptype{none} % do not increment counter
\begin{longtable}[]{@{}
  >{\raggedright\arraybackslash}p{(\linewidth - 4\tabcolsep) * \real{0.3333}}
  >{\raggedright\arraybackslash}p{(\linewidth - 4\tabcolsep) * \real{0.3333}}
  >{\raggedright\arraybackslash}p{(\linewidth - 4\tabcolsep) * \real{0.3333}}@{}}
\toprule\noalign{}
\begin{minipage}[b]{\linewidth}\raggedright
Quantity
\end{minipage} & \begin{minipage}[b]{\linewidth}\raggedright
Behavior
\end{minipage} & \begin{minipage}[b]{\linewidth}\raggedright
Source
\end{minipage} \\
\midrule\noalign{}
\endhead
\bottomrule\noalign{}
\endlastfoot
Number of relational waves \(M\) & \(O(N^2)\), unbounded & Definition \\
Number of terms of a decomposition \(\mathcal N\) (layer of existence) &
\textbf{Does not saturate} (the axioms admit any \(\mathcal N\)) & Axiom
1 \\
Number of bodies \(N\) (model specification) & External input
(structural specification of the representation) & Section 4.1 \\
Generator rank & Capped at \(\le 2N\) & Paper 3 {[}2{]} \\
Uniquely readable spatial directions & Saturates at 3 & Paper 3
{[}2{]} \\
Readout registers & Frozen under a fixed generator (no-go) & Paper 4
{[}3{]} \\
\end{longtable}
}

\[
\boxed{\ \text{Existence opens without limit; only the readout saturates}\ }
\]

A world that looks finite is a consequence of readout saturation, not of
the finiteness of existence. All ``splitting'' in this paper is a
statement of the readout layer --- register reallocation and the waxing
and waning of active planes --- and the operation of increasing the term
count \(\mathcal N\) is not an event requiring dynamical justification
(see limitation (i) in Section 9).

\subsubsection{1.3 Operational definition of
spontaneity}\label{operational-definition-of-spontaneity}

Applying a splitting map by hand proves no spontaneity. In this paper
spontaneity is defined operationally as dynamical instability.

\begin{enumerate}
\def\labelenumi{\arabic{enumi}.}
\tightlist
\item
  The growth rate of the dormant component does not depend on the seed
  amplitude \(\delta\) (exponential growth at an intrinsic rate).
\item
  Absolute-threshold crossing times are linear in \(\log(1/\delta)\);
  relative-threshold crossing times are \(\delta\)-independent.
\item
  The closure is conserved exactly throughout (the splitting is not
  caused by a breaking of closure).
\end{enumerate}

As the strongest verification we place \(\delta=0\). The applied
perturbation is zero; under exact arithmetic the orbit is invariant. If
the only seed is the nonzero numerical perturbation of floating-point
rounding, and that seed is amplified into collapse, this is strong
evidence of transverse instability --- the parent state collapses under
arbitrarily small perturbations.

\subsubsection{1.4 What this paper
decides}\label{what-this-paper-decides}

\begin{enumerate}
\def\labelenumi{\arabic{enumi}.}
\tightlist
\item
  How is the totality of closure-preserving dynamics characterized?
  Where does the fixed generator sit inside it?
\item
  Is the minimal intrinsic extension (unfreezing) consistent with Axiom
  0.5?
\item
  Does a self-consistent parent state (an exact periodic orbit of the
  unfrozen dynamics) exist? What is its structure?
\item
  Is the parent state unstable? Do the three criteria of spontaneity
  hold?
\item
  Can the outcomes after onset be classified? What decides class
  selection?
\item
  Is all of the above scale-covariant?
\end{enumerate}

No background space, distance matrix, or external observer is
introduced. No cosmological or field-theoretic physical names are used
in this paper (Axiom 0).

\begin{center}\rule{0.5\linewidth}{0.5pt}\end{center}

\subsection{2. Classification of Claims}\label{classification-of-claims}

{\def\LTcaptype{none} % do not increment counter
\begin{longtable}[]{@{}
  >{\raggedright\arraybackslash}p{(\linewidth - 4\tabcolsep) * \real{0.3333}}
  >{\raggedright\arraybackslash}p{(\linewidth - 4\tabcolsep) * \real{0.3333}}
  >{\raggedright\arraybackslash}p{(\linewidth - 4\tabcolsep) * \real{0.3333}}@{}}
\toprule\noalign{}
\begin{minipage}[b]{\linewidth}\raggedright
Subject
\end{minipage} & \begin{minipage}[b]{\linewidth}\raggedright
Classification
\end{minipage} & \begin{minipage}[b]{\linewidth}\raggedright
Position in this paper
\end{minipage} \\
\midrule\noalign{}
\endhead
\bottomrule\noalign{}
\endlastfoot
N-body complete pairwise relational-wave model & Model definition &
Implementation of Axiom 16 {[}1{]}; inherited from {[}2,3{]} \\
Non-constraint of \(\mathcal N\) (the term count is not internally
determined and does not saturate) & Consequence of Axiom 1 & Section
1.2. No ceiling in the layer of existence \\
The measurement target of this paper is the readout layer (not the layer
of existence) & Explicit framework & Section 1.2. All finiteness is
readout saturation \\
Scale anonymity & Axiom (first appearance) & Axiom 0.5. Unprovable
formation rule. Section 3 \\
The choice of zero on the right-hand side of Axiom 1 & Consequence of
Axiom 0.5 & The unique scale-free quadratic closure. Section 3 \\
Exact conservation of the complex square sum under orthogonal updates &
Mathematical consequence & Theorem 1 \\
Exhaustion of closure-preserving dynamics by \(\dot X=K(X)X\) &
Mathematical consequence & Theorem 2. Closure of the search space \\
Zeroth-order homogeneity of the unfrozen sine generator & Mathematical
consequence & Theorem 3. First application of Axiom 0.5 \\
Vanishing of the generator on collinear-phase states & Mathematical
consequence & Theorem 4. The parent must be circularly polarized \\
Exact periodicity of the eigenmode parent & Mathematical consequence &
Proposition 1 \\
Effective \(\operatorname{rank}2\) of the self-consistent parent's
generator (\(N=4\)) & Numerical insight & Unproven. Section 5.
Multi-plane for larger \(N\) (§5.3) \\
Exponential instability of the parent (rate independent of \(\delta\)) &
Numerical confirmation & Experiment O1. Section 8 \\
Collapse from rounding noise at \(\delta=0\) & Numerical confirmation &
Experiment O1. Section 8 \\
No-go under the fixed generator & Derived consequence {[}1,3{]} &
Numerically reconfirmed in control Experiment N \\
Threefold outcome classification (expansion, bounded mixing, recurrence)
& Numerical classification & Experiment O4. Classification thresholds
preliminary \\
Impossibility of expansion at \(N=3\) & Derived consequence & A
skew-symmetric matrix with \(M=3\) has \(\operatorname{rank}\le2\) \\
\(N\) phase structure (\(N\le4\) no expansion; \(N\ge5\) expansion in
all trials) & Numerical confirmation & Experiment O4 \\
Outcome stratification by the conserved ratio
\(c=\lvert Z^{\mathsf T}Z\rvert/\lVert Z\rVert^2\) (expansion suppressed
for \(N\le4\) at \(c=0\)) & Numerical confirmation & Experiment O5
(controlled contrast). Section 8.7 \\
Persistence of the onset law at large \(N\) (\(N=40\)--\(1000\), seed
depth \(f(0)\approx10^{-30}\), constant-rate amplification over 20+
decades) & Numerical confirmation & Experiment O6. Section 8.8 \\
Spontaneous exponential decay of \(\sigma_2\) (refreezing candidate) &
Numerical observation / working hypothesis & Confirmation of refreezing
requires additional criteria such as
\(\lVert K(\tau{+}1)-K(\tau)\rVert\). Sequel \\
Origin of the saturation ceiling \(f\approx0.335\)--\(0.338\) &
Unresolved & No identification with any specific constant is made \\
Relation to modulational instability {[}5,6{]} & Structural similarity &
Isomorphic phenomenology, different coupling structure (Section 9) \\
Recurrence of splitting and recombination & Numerical confirmation &
Phenomenology isomorphic to FPUT-type recurrence {[}7,8{]} \\
An explicit splitting map increasing the component count & Not
constructed & Measured here as growth of the active-plane count.
Connection problem \\
Stopping mechanism of splitting & Connection problem & Sequel (second
question) \\
\end{longtable}
}

\begin{center}\rule{0.5\linewidth}{0.5pt}\end{center}

\subsection{3. Axiom 0.5: Scale
Anonymity}\label{axiom-0.5-scale-anonymity}

\subsubsection{3.1 Full statement of the
axiom}\label{full-statement-of-the-axiom}

\begin{quote}
\textbf{Axiom 0.5 (scale anonymity)} The relations, update rules, and
readouts of the system may reference no absolute scale. Every
construction must be covariant under \(Z\to\lambda Z\)
(\(\lambda\in\mathbb C^\ast\)): the relations (closure conditions) are
homogeneous, the generator is zeroth-order homogeneous
(\(K(\lambda Z)=K(Z)\)), the state map is first-order homogeneous, and
the readouts are ratio quantities only. No means exists to measure an
absolute scale from inside the system.
\end{quote}

\subsubsection{3.2 Position}\label{position}

Axiom 0 is the formation rule ``no physical names may be referenced.''
Axiom 0.5 is the coequal formation rule ``no absolute scale may be
referenced,'' placed immediately after Axiom 0 and before Axiom 1. The
two are two clauses of a single idea: no names, no scale --- only
relations and ratios can be read.

This axiom is not a derived consequence. It is unprovable in the sense
that all readouts are ratios and no means of measuring an absolute scale
from inside the system exists in principle. It has the same rank as
Axiom 1 (zero-square-sum closure), which is likewise unprovable.

\subsubsection{3.3 The choice of the form of Axiom
1}\label{the-choice-of-the-form-of-axiom-1}

As a quadratic closure condition, any of \(\sum_n x_n^2=c\) was
conceivable; but \(c\ne0\) breaks under \(Z\to\lambda Z\) as
\(c\to\lambda^2c\). Within the family \(\sum_n x_n^2=c\), the only
dilation-invariant right-hand side is

\[
\sum_n x_n^2=0
\]

(\(c=0\)). If Axiom 0.5 is placed first, the zero on the right-hand side
of Axiom 1 is fixed with no room for choice within this family. If
general quadratic forms \(Z^{\mathsf T}GZ=0\) are admitted, other
homogeneous closures can be written, so the claim is not ``the quadratic
closure itself is unique'' but ``zero is the unique right-hand side of
\(\sum x_n^2=c\).'' The choice of the form \(\sum x_n^2\) (\(G=I\)) is
required by the name-permutation symmetry (anonymity). Moreover, the
compact part \(|\lambda|=1\) of the covariance group \(\mathbb C^\ast\)
is the common-phase invariance of the zero-square-sum closure: the
unreadability of phase and the unreadability of scale are the argument
part and the modulus part of one and the same group.

\subsubsection{3.4 Consequences within this
paper}\label{consequences-within-this-paper}

Axiom 0.5 performs three concrete jobs in this paper. First, it cuts the
search space of extended update rules: the generator must be
zeroth-order homogeneous and cannot depend on absolute amplitudes
(phases and amplitude ratios only). Second, it restricts all readout
thresholds to relative values: the activity threshold for planes, the
seed amplitude, and the recurrence criterion are all defined as ratios.
Third, it provides the acceptance criterion for numerical verification:
the ratio observables of a scale-transformed orbit must agree exactly
(Section 8.4).

After publication of this paper, this axiom will be incorporated into
layer 0 in the next version (v8) of the basic axiom system; Axiom 0 and
Axiom 0+ will then be merged as clauses of a single Axiom 0 (anonymity).

\begin{center}\rule{0.5\linewidth}{0.5pt}\end{center}

\subsection{4. Preliminaries: Model and Conservation
Structure}\label{preliminaries-model-and-conservation-structure}

\subsubsection{4.1 State and generator}\label{state-and-generator}

Consider the complete pairwise relational waves of \(N\) bodies. The
relational waves correspond to the edges of the complete graph \(K_N\);
there are \(M=\binom N2\) of them. The state is a complex vector
\(Z\in\mathbb C^M\) with phases \(\theta_e=\arg Z_e\). The line-graph
adjacency matrix \(A\) (\(A_{ef}=1\) when edges \(e,f\) share a body)
and the phase-difference sine generator

\[
K_{ef}=A_{ef}\sin(\theta_f-\theta_e)
\]

are inherited from the definitions of Paper 3 {[}2{]} (since
\(\theta_e=\arg Z_e\) is undefined at \(Z_e=0\), we adopt the convention
\(\arg 0:=0\); in the initial conditions of this paper all components
are nonzero). \(K\) is automatically skew-symmetric and is normalized in
spectral norm (after normalization \(\sigma_{\max}=1\), which makes all
subsequent thresholds relative). The update is the Cayley transform
{[}10{]}

\[
U=(I-\gamma K)^{-1}(I+\gamma K),
\qquad
\gamma=\tan\frac{\pi}{144}.
\]

If \(K\) is skew-symmetric, \(U\) is a real orthogonal matrix.

\textbf{Fixed generator (control)}: \(K\) is evaluated at the phases
\(\arg Z(0)\) at \(\tau=0\) and frozen. This is the dynamics of Papers 3
and 4.

\textbf{Unfreezing (the minimal extension of this paper)}: \(K\) is
reconstructed at every step from the current phases \(\arg Z(\tau)\).
What is introduced is not a new coupling rule but only the release of
the freezing of a generator that was defined as a function of phase
differences from the start.

\subsubsection{4.2 Theorem 1 (exact conservation of the complex square
sum under orthogonal
updates)}\label{theorem-1-exact-conservation-of-the-complex-square-sum-under-orthogonal-updates}

An update \(Z'=UZ\) by a real orthogonal matrix \(U\) conserves the
complex bilinear square sum and the Hermitian norm simultaneously and
exactly:

\[
Z'^{\mathsf T}Z'=Z^{\mathsf T}U^{\mathsf T}UZ=Z^{\mathsf T}Z,
\qquad
Z'^{\dagger}Z'=Z^{\dagger}Z.
\]

\textbf{Proof} Immediate from \(U^{\mathsf T}U=I\) and the realness of
\(U\). \(\blacksquare\)

Conservation of the complex square sum \(Z^{\mathsf T}Z\) is the
simultaneous conservation of two conditions, its real and imaginary
parts. These are the same two conditions demanded by the local splitting
condition \(X_P^2=\sum_kY_k^2\); the dynamics of this paper never breaks
this closure at any point of the splitting process.

\subsubsection{4.3 Theorem 2 (exhaustion of closure-preserving
dynamics)}\label{theorem-2-exhaustion-of-closure-preserving-dynamics}

Every smooth flow on the closure sphere \(\{X:\,X^{\mathsf T}X=R^2\}\)
(\(R^2>0\), \(X\ne0\)) can be written, with a skew-symmetric
matrix-valued function \(K(X)\), in the form

\[
\dot X=K(X)X,
\]

and conversely every flow of this form preserves the closure.

\textbf{Proof} Conservation:
\(\frac{d}{d\tau}X^{\mathsf T}X=2X^{\mathsf T}K(X)X=0\) (skew-symmetry).
Exhaustion: for a tangent vector field \(V(X)\) on the sphere
(\(X^{\mathsf T}V=0\)), set
\(K(X)=\dfrac{VX^{\mathsf T}-XV^{\mathsf T}}{X^{\mathsf T}X}\); then
\(K\) is skew-symmetric and \(K(X)X=V\). \(\blacksquare\)

\textbf{Remark (scope)} (i) For isotropic states (\(X^{\mathsf T}X=0\),
\(X\ne0\)) the denominator of the above construction vanishes; but
taking \(Y\) with \(Y^{\mathsf T}X=1\) and setting
\(K=VY^{\mathsf T}-YV^{\mathsf T}\) gives \(KX=V-Y(V^{\mathsf T}X)=V\)
(using the tangency condition \(V^{\mathsf T}X=0\)), which repairs the
construction. \(X=0\) is excluded. (ii) The theorem exhausts
continuous-time flows; the state-dependent Cayley update of this paper
is one discretization within this class. No exhaustion of all discrete
maps is claimed.

Thus a ``state-dependent generator'' is not one of many extension
candidates but the totality of closure-preserving dynamics. The fixed
generator is merely the special case \(K(X)=K_0\), and the search for
the minimal extension enabling splitting may be conducted inside this
space.

\subsubsection{4.4 Theorem 3 (scale covariance of the unfrozen
generator)}\label{theorem-3-scale-covariance-of-the-unfrozen-generator}

The unfrozen sine generator is zeroth-order homogeneous: for any
\(\lambda\in\mathbb C^\ast\),

\[
K(\lambda Z)=K(Z).
\]

\textbf{Proof} Write \(\lambda=|\lambda|e^{i\varphi}\). Then
\(\arg(\lambda Z_e)=\arg Z_e+\varphi\), and the phase differences
\(\theta_f-\theta_e\) depend on neither \(|\lambda|\) nor \(\varphi\).
Spectral-norm normalization is a ratio and hence invariant.
\(\blacksquare\)

The unfrozen dynamics therefore satisfies Axiom 0.5 without adjustment.
Since the update is linear (\(U\) is a zeroth-order function of \(Z\)),
the entire orbit commutes with \(Z\to\lambda Z\), and the orbits of
ratio observables are exactly scale-independent.

\subsubsection{\texorpdfstring{4.5 Two invariants and the conserved
ratio
\(c\)}{4.5 Two invariants and the conserved ratio c}}\label{two-invariants-and-the-conserved-ratio-c}

By Theorem 1 the dynamics conserves two quantities simultaneously and
exactly: the Hermitian norm \(\lVert Z\rVert^2\) and the complex
bilinear square sum \(Z^{\mathsf T}Z\). Their ratio

\[
c=\frac{\lvert Z^{\mathsf T}Z\rvert}{\lVert Z\rVert^2}\in[0,1]
\]

is a conserved dimensionless ratio and a readable quantity conforming to
Axiom 0.5 (the upper bound \(c\le1\) follows from the triangle
inequality \(\lvert\sum Z_e^2\rvert\le\sum\lvert Z_e\rvert^2\); equality
is equivalent to squared-phase alignment of all components). \(c\)
stratifies the orbits: states with different \(c\) cannot be carried
into one another by the dynamics.

\begin{itemize}
\tightlist
\item
  \textbf{\(c=0\) stratum}: zero closure \(Z^{\mathsf T}Z=0\). The
  stratum designated by Axiom 1; Experiments O1 and O4 of this paper run
  here
\item
  \textbf{\(c>0\) strata}: the closure convention
  \(X^{\mathsf T}X=R^2>0\) of Papers 3 and 4 {[}2,3{]} corresponds here.
  The theorems of both papers are linear algebra and do not depend on
  \(c\), but the standard initial states of their experiments (imaginary
  seed amplitude \(s=0.35\)) run in the intermediate stratum
  \(c=R^2/(R^2+2s^2)\approx0.80\). Purely real states (for which
  \(X^{\mathsf T}X=\lVert X\rVert^2\) holds identically) are the
  limiting case \(c=1\)
\item
  Generic complex states take intermediate \(c\)
\end{itemize}

That the series has moved back and forth between ``real sphere'' and
``zero closure'' formulations is not a contradiction but a use of strata
labeled by the conserved quantity \(c\). And as Section 8.7 shows, at
small \(N\) the value of \(c\) actually changes the outcome class ---
the distinction of strata is physics, not representational convenience.

\begin{center}\rule{0.5\linewidth}{0.5pt}\end{center}

\subsection{5. The Parent State: Necessity of the Circularly Polarized
Eigenmode}\label{the-parent-state-necessity-of-the-circularly-polarized-eigenmode}

\subsubsection{5.1 Theorem 4 (vanishing of the generator on
collinear-phase
states)}\label{theorem-4-vanishing-of-the-generator-on-collinear-phase-states}

On a state whose component phases are collinear
(\(\theta_e\in\{\psi,\psi+\pi\}\) for some \(\psi\)), the sine generator
is exactly zero.

\textbf{Proof} For any \(e,f\) we have
\(\theta_f-\theta_e\in\{0,\pm\pi\}\), hence
\(\sin(\theta_f-\theta_e)=0\). \(\blacksquare\)

A state combining plane basis vectors with real coefficients,
\(Z=e^{i\psi}(\cos\varphi\,p+\sin\varphi\,q)\) (\(p,q\) real vectors),
falls under this condition since every component is a real multiple of
\(e^{i\psi}\). That is, linearly polarized in-plane states degenerate
the generator and cannot carry self-consistent dynamics. \textbf{The
exact parent constructed in this paper is the circularly polarized
eigenmode.} What Theorem 4 excludes is collinear-phase states;
Proposition 1 gives a sufficient condition for a circularly polarized
eigenmode to be an exact periodic parent. That no other kind of
non-collinear self-consistent periodic orbit exists is not shown (open).

\subsubsection{5.2 Proposition 1 (exact periodicity of the eigenmode
parent)}\label{proposition-1-exact-periodicity-of-the-eigenmode-parent}

Let \((\sigma,p,q)\) be the maximal singular-value plane of \(K\)
(orientation convention \(Kp=+\sigma q\)) and

\[
v=\frac{p+iq}{\sqrt2},
\qquad
Kv=-i\sigma v
\]

its eigenmode. If \(v\) is self-consistent with its own generator
(\(K(\arg v)=K\)), then under the unfrozen dynamics \(v\) is an exact
periodic orbit.

\textbf{Proof}
\(Uv=\dfrac{1-i\gamma\sigma}{1+i\gamma\sigma}\,v=e^{-i\theta}v\)
(\(\theta=2\arctan\gamma\sigma\)). The update is a common phase
rotation; the phases of all components shift uniformly, so the phase
differences are invariant. Hence \(K(\arg Uv)=K(\arg v)\) is preserved,
and the same rotation repeats at every subsequent step. \(\blacksquare\)

The self-consistent eigenmode was constructed numerically by the
fixed-point iteration \(\theta\leftarrow\arg v(K(\theta))\). Since the
in-plane \(SO(2)\) gauge acts only on the common phase of \(v\)
(rotating \(p,q\) by an angle \(\alpha\) sends \(v\to e^{-i\alpha}v\)),
the iteration does not depend on the gauge choice. The reached parent
had a median eigenmode-equation residual of \(8.9\times10^{-13}\); in
the \(\delta=0\) verification, over one period (144 steps) the generator
variation stayed at \(9.2\times10^{-14}\) and the dormant fraction
remained exactly zero.

\subsubsection{\texorpdfstring{5.3 Numerical insight: the parent's
generator has
\(\operatorname{rank}2\)}{5.3 Numerical insight: the parent's generator has \textbackslash operatorname\{rank\}2}}\label{numerical-insight-the-parents-generator-has-operatornamerank2}

The generator \(K(\arg v)\) of the self-consistent parent had effective
\(\operatorname{rank}2\) (one rotation plane, \(\sigma=1\)) in every
configuration examined at \(N=4\), the remaining \(M-2\) dimensions
being the kernel (static subspace). The parent state sees only its own
plane; all dormant directions belong to the kernel. No proof of this
property has been obtained, and it is likely specific to \(N=4\) ---
separately constructed self-consistent parents at larger \(N\) had
generators with several planes
(\(\sigma_2/\sigma_1\approx0.35\)--\(0.5\); treated in the sequel). The
design of the dormant seed in this paper (kernel directions) is based on
this \(N=4\) structure.

The dormant seed is placed so as to preserve zero closure: constructing
\(g=(u+iw)/\sqrt2\) from a real orthonormal pair \((u,w)\) in the kernel
gives \(g^{\mathsf T}g=0\) exactly, and together with
\(v^{\mathsf T}g=0\) the initial state preserves zero closure. For the
\(N\ge5\) parent families of Sections 8.6 and 8.7 the parent generator
is multi-plane and the ``kernel direction'' design does not apply; the
seeds of those families are constructed by projection onto the
orthogonal complement of the parent plane \(\mathrm{span}\{v,\bar v\}\)
(Experiment O4: complex Gaussian projection; Experiment O5 parentZ0: a
zero-closure pair from a real orthonormal pair inside the complement).
At \(N=3\) the orthogonal complement of the parent plane is only
one-dimensional and no zero-closure perturbation orthogonal to the
parent exists (an isotropic vector needs two real dimensions) --- in
that case a real seed \(g=u\) is used and a closure residual of
\(O(\delta^2)\) is explicitly tolerated (the \(N=3\) parent family of
Section 8.6).

In addition, in 4 of 6 parent preparations at \(N=5\) and 1 of 6 at
\(N=6\), the eigenmode residual stalled at \(\sim9\times10^{-2}\)
(worsened convergence of the fixed-point iteration, or possibly
nonexistence of the fixed point at larger \(N\)). The ``parents'' of
those trials are approximate parents. The conclusions of Section 8.6 are
unchanged when restricted to trials with residuals \(\sim10^{-15}\) and
are robust, but the existence problem of parents at larger \(N\) remains
open.

\begin{center}\rule{0.5\linewidth}{0.5pt}\end{center}

\subsection{6. Readout Definitions for
Splitting}\label{readout-definitions-for-splitting}

Splitting is read through three ratio quantities (all relative
quantities conforming to Axiom 0.5).

\textbf{Dormant fraction} For the parent plane \((p_1,q_1)\) of the
reference generator \(K(\arg Z(0))\),

\[
f(\tau)=1-\frac{|p_1^{\mathsf T}Z|^2+|q_1^{\mathsf T}Z|^2}{Z^{\dagger}Z}.
\]

The share of the allocation outside the parent plane (in the kernel).
This is the readout of splitting as register reallocation.

\textbf{Second singular direction} The second singular value
\(\sigma_2\) of the instantaneous generator \(K(\arg Z(\tau))\) (after
normalization \(\sigma_1=1\), so \(\sigma_2\) is a ratio). The relative
threshold \(\sigma_2>0.05\) defines an active plane. This is the readout
of the establishment of a new rotation mode --- corresponding to an
increase of the component count. Note that an absolute threshold cannot
be placed by Axiom 0.5; numerical dust \(\sigma_2=O(\delta)\) must not
be counted as a plane.

\textbf{Participation ratio} The participation ratio of edge amplitudes

\[
\mathrm{PR}(\tau)=\frac{\left(\sum_e|Z_e|^2\right)^2}{\sum_e|Z_e|^4}\in[1,M].
\]

The readout of the spread of the allocation {[}11{]}.

\begin{center}\rule{0.5\linewidth}{0.5pt}\end{center}

\subsection{7. Experimental Design}\label{experimental-design}

\subsubsection{Experiment N (control)}\label{experiment-n-control}

The identical initial condition (parent + dormant seed) is updated with
the fixed generator (frozen at \(\tau=0\)). By Consequences 16.5 and
16.6 {[}1,3{]}, the reference plane registers and the kernel allocation
should be individually conserved to machine precision. This doubles as a
soundness check of the measurement system and a numerical reconfirmation
of the no-go.

\subsubsection{Experiment O1 (onset of
splitting)}\label{experiment-o1-onset-of-splitting}

The identical initial condition is updated with unfreezing. \(N=4\)
(\(M=6\)), seed amplitudes \(\delta=10^{-2},10^{-3},10^{-4},10^{-5}\), 6
random seeds each, 4320 steps. The dormant seed is placed from a real
orthonormal pair in the kernel as \(g=(u+iw)/\sqrt2\)
(\(g^{\mathsf T}g=0\) exact), and the initial state satisfies zero
closure to \(|Z^{\mathsf T}Z|\le3.2\times10^{-13}\). The initial dormant
fraction came to \(f(0)\approx0.62\,\delta^2\). In addition a
\(\delta=0\) verification run (pure eigenmode, no perturbation) is
placed.

\subsubsection{Scale verification}\label{scale-verification}

For a representative condition (\(\delta=10^{-3}\)), the ratio
observables of orbits with \(Z_0\to\lambda Z_0\) are compared:
\(\lambda=2^{10}=1024\) (a scale change with exactly zero floating-point
rounding error), \(\lambda=987.654\) (generic value),
\(\lambda=e^{0.7i}\) (common phase rotation). Additionally the
orbit-divergence rate against a twin perturbation (one component changed
by a relative \(10^{-13}\)) is compared, to separate deviations at
generic \(\lambda\) --- dynamical amplification of rounding differences
--- from any breaking of scale symmetry.

\subsubsection{Experiment O4 (outcome
classification)}\label{experiment-o4-outcome-classification}

\(N=3,4,5,6\) (\(M=3,6,10,15\)), two families of initial conditions
(parent: self-consistent parent + \(\delta=10^{-3}\) seed; generic:
uniformly random complex state), 6 random seeds each, 20000 steps. The
classification metrics are the late-window (second-half) minimum of
\(f\) (recurrence depth), the residence fraction of \(\sigma_2>0.05\)
(establishment of expansion), and the participation ratio. The
(preliminary) classification rules are:

\begin{itemize}
\tightlist
\item
  \textbf{Expansion}: \(\sigma_2\) active residence fraction \(>0.5\)
\item
  \textbf{Recurrence}: late-window \(\min f<0.05\)
\item
  \textbf{Bounded mixing}: everything else
\end{itemize}

\subsubsection{Experiment O5 (controlled contrast of the closure
convention)}\label{experiment-o5-controlled-contrast-of-the-closure-convention}

With everything else identical (dynamics, classification metrics, 20000
steps), four families contrast only the initial conserved ratio \(c\):
genericZ0 (generic state, \(Z^{\mathsf T}Z=0\) exact), genericNZ
(generic state, complex Gaussian, \(c\approx0.3\)--\(0.6\)), parentZ0
(parent + zero-closure kernel seed), parentNZ (parent + complex Gaussian
seed, \(c\lesssim\delta^2\)). \(N=3\)--\(6\), 3 random seeds each.
Whether \(c\) changes the outcome class is judged directly.

\subsubsection{Experiment O6 (large-N onset
runs)}\label{experiment-o6-large-n-onset-runs}

To confirm that the onset law (Experiment O1) is not a peculiarity of
small \(N\), onset runs at \(N=40\) (\(M=780\)), \(N=300\)
(\(M=44850\)), and \(N=1000\) (\(M=499500\)) are performed with the
low-rank vertex-decomposition engine. The engine uses the vertex
decomposition \(K=WJW^{\mathsf T}\) of Paper 3 as an algorithm, reducing
one step to \(O(N^3)\) (agreement with the dense implementation verified
at \(N=12\): maximum orbit deviation \(5.6\times10^{-14}\) over 300
steps).

The seed is \(f(0)\approx\delta^2=10^{-30}\) in energy ratio
(\(\delta=10^{-15}\), near the lower limit at which the seed's relative
amplitude against the parent components approaches the double-precision
rounding unit). To measure this depth, the dormant fraction is measured
not by the subtraction \(1-h_1/h_{\mathrm{total}}\) (whose cancellation
floor is \(10^{-16}\)) but directly by the projection norm onto the
orthogonal complement of the parent plane,
\(f=\lVert Z-p(p^{\mathsf T}Z)-q(q^{\mathsf T}Z)\rVert^2/\lVert Z\rVert^2\).
The parent's self-consistent iteration is converged to residuals at the
\(10^{-12}\)--\(10^{-11}\) level (achieved residuals are listed in the
results table; the energy of the construction residual,
\(\sim10^{-24}\)--\(10^{-23}\), sets the measurement floor). The seed is
\(g=(u+iw)/\sqrt2\) (\(g^{\mathsf T}g=0\) exact) from a real orthonormal
pair in the kernel (for large \(N\), the orthogonal complement of
\(\operatorname{span}(W)\)). One random seed per \(N\); runs are tracked
to 1500 steps past the crossing of the threshold \(f>0.05\) (cap 9000
steps).

\begin{center}\rule{0.5\linewidth}{0.5pt}\end{center}

\subsection{8. Results}\label{results}

\subsubsection{8.1 Control (Experiment N): numerical reconfirmation of
the
no-go}\label{control-experiment-n-numerical-reconfirmation-of-the-no-go}

Under the fixed generator, the maximum variation of the reference
register ratios was \(\le4.5\times10^{-11}\) across all trials, and the
dormant kernel allocation remained at \(\delta^2\) without any growth
--- exactly as Consequences 16.5 and 16.6 predict.

\begin{figure}
\centering
\pandocbounded{\includegraphics[keepaspectratio,alt={Fixed generator versus unfreezing}]{figs/register_migration_fixed_vs_sequential_v1.png}}
\caption{Fixed generator versus unfreezing}
\end{figure}

\emph{Figure 1: Orbits of the reference plane registers (\(N=4\),
\(\delta=10^{-3}\)). Left: fixed generator (control) --- the parent
plane and kernel allocations are frozen to machine precision. Right:
unfreezing --- from the identical initial condition the allocation moves
spontaneously.}

\subsubsection{8.2 Onset of splitting (Experiment O1): the three
criteria of spontaneity
hold}\label{onset-of-splitting-experiment-o1-the-three-criteria-of-spontaneity-hold}

\[
\begin{array}{c|cccc}
\delta & 10^{-2} & 10^{-3} & 10^{-4} & 10^{-5}\\\hline
\text{initial growth rate (median, /step}) & 3.795\times10^{-2} & 3.783\times10^{-2} & 3.781\times10^{-2} & 3.781\times10^{-2}\\
\tau_c(f\ge10^{-2})\text{ (median}) & 137.5 & 222 & 306 & 390\\
\tau_c(f\ge10f(0))\text{ (median}) & 83.5 & 84 & 84 & 84
\end{array}
\]

The growth rate is invariant over three decades of \(\delta\) (the rate
depends on the random seed --- the parent configuration --- but not on
\(\delta\)). Absolute-threshold crossing times are linear in
\(\log(1/\delta)\) (about 84 steps per decade); relative-threshold
crossing times are \(\delta\)-independent. Criteria 1 and 2 hold.

The growth is not single-exponential, however. Against the crossing
interval \(2\ln10/g\approx122\) steps predicted from the
initial-fit-window rate \(g\approx3.8\times10^{-2}\), the measured
interval is about 84 steps, so the effective rate defined from crossing
times is \(2\ln10/84\approx5.5\times10^{-2}\). The initial transient and
the main regime have different rates; both values are reported as a
limitation of the single-exponential approximation.

\begin{figure}
\centering
\pandocbounded{\includegraphics[keepaspectratio,alt={Dormant-fraction growth for the δ series}]{figs/dormant_fraction_growth_delta_series_v1.png}}
\caption{Dormant-fraction growth for the δ series}
\end{figure}

\emph{Figure 2: Growth of the dormant fraction \(f(\tau)\) (median over
random seeds). Four parallel exponential ramps shifted by
\(\log(1/\delta)\) --- the classic signature of an instability. After
saturation the curves enter a universal ceiling
(\(f\approx0.335\)--\(0.338\)) and back-and-forth oscillation.}

\begin{figure}
\centering
\pandocbounded{\includegraphics[keepaspectratio,alt={δ dependence of crossing times}]{figs/crossing_time_vs_delta_v1.png}}
\caption{δ dependence of crossing times}
\end{figure}

\emph{Figure 3: Threshold crossing times. Absolute thresholds (circles,
squares) are linear in \(\log_{10}(1/\delta)\); the relative threshold
(triangles) is \(\delta\)-independent.}

\begin{figure}
\centering
\pandocbounded{\includegraphics[keepaspectratio,alt={δ independence of the growth rate}]{figs/growth_rate_vs_delta_v1.png}}
\caption{δ independence of the growth rate}
\end{figure}

\emph{Figure 4: Distribution of initial exponential growth rates. The
point pattern of each \(\delta\) column is identical --- the rate is
decided by the random seed (parent configuration), not by the seed
amplitude.}

Decisive is \(\delta=0\). The applied perturbation is zero, and under
exact arithmetic the orbit would be invariant. Nevertheless the orbit
collapsed from the rounding-noise floor (\(f\sim10^{-33}\)) through an
exponential ramp of about 20 decades, reaching \(f>10^{-10}\) at
\(\tau\approx850\) and then saturating to the same ceiling ---
amplification of the nonzero numerical perturbation of rounding, and
strong evidence of transverse instability.

\begin{figure}
\centering
\pandocbounded{\includegraphics[keepaspectratio,alt={Collapse from rounding noise at δ=0}]{figs/zero_delta_machine_noise_onset_v1.png}}
\caption{Collapse from rounding noise at δ=0}
\end{figure}

\emph{Figure 5: The \(\delta=0\) verification. An orbit that was exactly
periodic over one period (generator variation \(9.2\times10^{-14}\),
\(f=0\)) collapses spontaneously with nothing but floating-point
rounding as the seed.}

\subsubsection{8.3 Exact conservation of the
closure}\label{exact-conservation-of-the-closure}

\[
\begin{array}{l|ccc}
 & |\Delta\operatorname{Re}Z^{\mathsf T}Z| & |\Delta\operatorname{Im}Z^{\mathsf T}Z| & |\Delta Z^{\dagger}Z|/Z^{\dagger}Z\\\hline
\text{fixed generator} & \le1.1\times10^{-14} & \le8.3\times10^{-15} & \le1.1\times10^{-12}\\
\text{unfreezing} & \le1.7\times10^{-14} & \le1.2\times10^{-14} & \le1.8\times10^{-14}
\end{array}
\]

The onset of splitting proceeds without breaking the real and imaginary
parts of the complex square-sum closure (the two conditions demanded by
the local splitting condition). Criterion 3 holds.

\subsubsection{8.4 Numerical verification of scale anonymity and
separation of
sensitivity}\label{numerical-verification-of-scale-anonymity-and-separation-of-sensitivity}

\[
\begin{array}{l|c}
\text{transformation} & \max|\Delta(\text{ratio observable})|\\\hline
\lambda=2^{10}=1024 & 0\ \text{(bitwise agreement)}\\
\lambda=987.654 & 2.9\times10^{-14}\\
\lambda=e^{0.7i}\ \text{(common phase)} & 1.1\times10^{-13}\\
\text{twin perturbation }10^{-13}\ \text{(reference)} & 1.6\times10^{-11}\\
\text{twin perturbation, fixed generator (reference)} & 7.8\times10^{-16}
\end{array}
\]

Bitwise agreement at \(\lambda=2^{10}\), where no rounding error occurs
at all, shows the implementation to be exactly scale-free. The
deviations \(\lesssim10^{-13}\) at generic \(\lambda\) have the same
slow amplification rate (\(\sim3.8\times10^{-4}\)/step) as the twin
perturbation and are dynamical amplification of initial rounding
differences, not a breaking of scale symmetry. Together with Theorem 3,
Axiom 0.5 holds both theoretically and numerically.

\begin{figure}
\centering
\pandocbounded{\includegraphics[keepaspectratio,alt={Closure conservation and sensitivity}]{figs/closure_conservation_and_sensitivity_v1.png}}
\caption{Closure conservation and sensitivity}
\end{figure}

\emph{Figure 6: Left --- conservation error of the complex square-sum
closure (machine precision for both dynamics). Right --- sensitivity
probes. Generic scale, common phase, and twin perturbation show
comparable slow divergence; the fixed-generator reference does not
diverge.}

\subsubsection{\texorpdfstring{8.5 After onset: saturation, recurrence,
non-establishment of a second plane
(\(N=4\))}{8.5 After onset: saturation, recurrence, non-establishment of a second plane (N=4)}}\label{after-onset-saturation-recurrence-non-establishment-of-a-second-plane-n4}

After saturation, \(f\) has a universal ceiling \(0.335\)--\(0.338\) and
oscillates between \(0.04\) and \(0.34\). This is repeated splitting and
recombination in a conservative system --- phenomenology isomorphic to
FPUT-type recurrence {[}7,8{]}. The origin of the ceiling value
(structural or dynamical constant) is unresolved; no identification with
any specific constant is made. Deciding this requires systematic
measurement of the \(N\) and \(\gamma\) dependence.

Meanwhile the instantaneous generator remains single-plane throughout.
\(\sigma_2\) stays at \(10^{-3}\)--\(10^{-4}\) and never exceeds the
relative threshold \(0.05\). Near the parent at \(N=4\), splitting as
register reallocation occurs, but the establishment of a new rotation
plane --- splitting as an increase of the component count --- does not.

\begin{figure}
\centering
\pandocbounded{\includegraphics[keepaspectratio,alt={Non-establishment of the second singular direction}]{figs/instantaneous_plane_count_v1.png}}
\caption{Non-establishment of the second singular direction}
\end{figure}

\emph{Figure 7: The second singular direction \(\sigma_2\) of the
instantaneous generator (\(N=4\), parent family). For all \(\delta\) it
stays below the activity threshold \(0.05\) --- no second rotation plane
establishes itself.}

\subsubsection{8.6 Threefold outcome classification (Experiment
O4)}\label{threefold-outcome-classification-experiment-o4}

\[
\begin{array}{cc|cc}
N & M & \text{parent family} & \text{generic family}\\\hline
3 & 3 & \text{bounded mixing }6/6 & \text{bounded mixing }5/6,\ \text{recurrence }1/6\\
4 & 6 & \text{bounded mixing }6/6\ (\sigma_2\text{ exp. decay}) & \text{bounded mixing }6/6\\
5 & 10 & \text{expansion }6/6 & \text{expansion }6/6\\
6 & 15 & \text{expansion }6/6 & \text{expansion }6/6
\end{array}
\]

Across all trials the maximum deviation of closure conservation was
\(3.3\times10^{-14}\), and the initial zero closure was
\(|Z^{\mathsf T}Z|\le2.3\times10^{-16}\) (except the \(N=3\) parent
family, where no zero-closure seed exists and \(O(\delta^2)=10^{-6}\) is
tolerated; Section 5.3). All three outcome classes were realized
(recurrence in one trial only).

\textbf{Expansion} A new rotation plane establishes itself
(\(\sigma_2>0.05\) residence fraction 1) and the allocation spreads over
all relational waves, converging to a stationary multi-plane state
(parent-family representatives: \(N=5\): \(f\to0.55\),
\(\sigma_2\approx0.25\) stationary; \(N=6\): \(f\to0.8\),
\(\sigma_2\approx0.4\) stationary).

\textbf{Bounded mixing} The allocation spreads (in the parent family the
participation ratio converges numerically to the equipartition value
\(\mathrm{PR}=M\); median deviation \(\lesssim10^{-2}\)), but no second
plane establishes itself; a single plane wanders while mixing amplitudes
only.

\textbf{Recurrence} Repeated deep returns toward the parent plane. In
the zero-closure stratum (\(c=0\)) it appeared in only one trial of the
\(N=3\) generic family (\(f_{\min}=0.021\)), but at \(c\approx0.55\),
\(N=3\) shows it systematically in 3/3 (Section 8.7).

\begin{figure}
\centering
\pandocbounded{\includegraphics[keepaspectratio,alt={Fate curves of outcomes}]{figs/fate_curves_parent_family_v1.png}}
\caption{Fate curves of outcomes}
\end{figure}

\emph{Figure 8: Fate curves of the parent family (\(N=3,4,5,6\)).
\(N=3\) has \(\sigma_2\equiv0\) (expansion structurally impossible). At
\(N=4\), \(\sigma_2\) decays spontaneously and exponentially from
\(6\times10^{-4}\) into the \(10^{-5}\) range. For \(N\ge5\),
\(\sigma_2\) settles at a stationary value and expansion completes.}

\begin{figure}
\centering
\pandocbounded{\includegraphics[keepaspectratio,alt={Outcome-class map}]{figs/outcome_class_map_v1.png}}
\caption{Outcome-class map}
\end{figure}

\emph{Figure 9: Classification map of all 48 trials (4 values of \(N\) ×
2 families × 6 random seeds; horizontal: late recurrence depth,
vertical: \(\sigma_2\) active residence fraction). Expansion (top) and
non-expansion (bottom) separate cleanly.}

\begin{figure}
\centering
\pandocbounded{\includegraphics[keepaspectratio,alt={Spread of the allocation}]{figs/participation_ratio_v1.png}}
\caption{Spread of the allocation}
\end{figure}

\emph{Figure 10: Late participation ratio. The parent family converges
numerically to \(\mathrm{PR}=M\) (equipartition) for all \(N\).}

The structure of class selection is summarized in three points.

\begin{enumerate}
\def\labelenumi{\arabic{enumi}.}
\tightlist
\item
  \textbf{\(N\) phase structure} At \(N=3\) expansion is structurally
  impossible: a skew-symmetric matrix with \(M=3\) has
  \(\operatorname{rank}\le2\), so no second rotation plane can exist
  (\(\sigma_2\equiv0\), derived). At \(N=4\) no expansion arose from
  zero-closure states either (both families bounded mixing). For
  \(N\ge5\) expansion arose in both families and all random seeds. The
  boundary lies between \(N=4\) and \(N=5\) --- though ``for all initial
  conditions examined'': this is not a theorem over all initial states.
\item
  \textbf{Outcome stratification by the conserved ratio \(c\)} The
  controlled contrast (Experiment O5, Section 8.7) established that what
  separates expansion at \(N=4\) is not proximity to the parent but the
  exactly conserved ratio \(c\). Restricting to \(c=0\) removes
  expansion for \(N\le4\). That is, the closure condition of Axiom 1
  (designation of the \(c=0\) stratum) is not only a conservation law
  but narrows the reachable end states of the dynamics. The mechanism is
  a connection problem.
\item
  \textbf{Observation of spontaneous refreezing} In the \(N=4\) parent
  family, \(\sigma_2\) decayed exponentially without external
  intervention (\(6\times10^{-4}\to10^{-5}\) over 20000 steps),
  suggesting the system heads spontaneously toward a new single-plane
  self-consistent configuration --- an effectively frozen state.
  However, \(\sigma_2\to0\) only means the instantaneous generator
  approaches \(\operatorname{rank}2\); if the orientation of the plane
  keeps moving, allocation transfer can continue. Confirming refreezing
  requires \(\lVert K(\tau{+}1)-K(\tau)\rVert_{\mathrm F}\to0\) or
  vanishing drift of the principal-plane projector, which is done in the
  sequel.
\end{enumerate}

\subsubsection{\texorpdfstring{8.7 Controlled contrast of the closure
convention (Experiment O5): the conserved ratio \(c\) stratifies
outcomes}{8.7 Controlled contrast of the closure convention (Experiment O5): the conserved ratio c stratifies outcomes}}\label{controlled-contrast-of-the-closure-convention-experiment-o5-the-conserved-ratio-c-stratifies-outcomes}

\[
\begin{array}{c|cccc}
N & \text{genericZ0 }(c\approx0) & \text{genericNZ }(c\approx0.3\text{–}0.6) & \text{parentZ0 }(c\approx0) & \text{parentNZ}\\\hline
3 & \text{bounded mixing }3/3 & \textbf{recurrence }3/3 & \text{bounded mixing }3/3 & \text{bounded mixing }3/3\\
4 & \text{bounded mixing }3/3 & \textbf{expansion }3/3 & \text{bounded mixing }3/3 & \text{bounded mixing }3/3\\
5 & \text{expansion }3/3 & \text{expansion }3/3 & \text{expansion }3/3 & \text{expansion }3/3\\
6 & \text{expansion }3/3 & \text{expansion }3/3 & \text{expansion }3/3 & \text{expansion }3/3
\end{array}
\]

Maximum deviation of closure conservation across all 48 runs:
\(4.8\times10^{-14}\). Three readings are established.

\begin{enumerate}
\def\labelenumi{\arabic{enumi}.}
\tightlist
\item
  \textbf{At small \(N\) (3, 4) the value of \(c\) changes the outcome
  class.} The controlled comparison between the generic families
  (everything identical, only \(c\) contrasted) separated systematically
  3/3. For \(N\ge5\) it has no effect (all expansion)
\item
  \textbf{What acts is the magnitude of \(c\).} The parentNZ family at
  \(N=3\) (\(c\approx10^{-6}\)) fell into the same class as the \(c=0\)
  side, not as genericNZ (\(c\approx0.55\)) --- not a binary ``zero
  versus nonzero'' but the magnitude of the conserved ratio decides the
  stratum
\item
  \textbf{The old ``near-parent versus generic'' family contrast was
  confounded with \(c\).} After separation by this controlled contrast,
  the variable separating the \(N=4\) outcome is not the character of
  the initial family but \(c\)
\end{enumerate}

The recurrence class gained systematic support from this contrast: at
\(c\approx0.55\), \(N=3\) recurs in 3/3; the recurrence that appeared in
only one trial in Experiment O4 (\(c=0\)) is a standard outcome in the
nonzero-\(c\) strata.

Reservation: in the parentNZ family, by the nature of the seed
construction, \(c\) came out unintentionally tiny (\(\approx10^{-16}\)
for \(N\ge4\), \(10^{-6}\) for \(N=3\)), so the parent-side contrast is
used only as evidence for point 2 (``tiny \(c\) behaves as \(c=0\)'').
The parentZ0 family at \(N=3\) degenerates (one-dimensional complement;
Section 5.3).

\subsubsection{8.8 Large-N onset runs (Experiment O6): body-count
universality of the onset law and seed depth of 30
decades}\label{large-n-onset-runs-experiment-o6-body-count-universality-of-the-onset-law-and-seed-depth-of-30-decades}

\[
\begin{array}{c|ccc}
 & N=40\ (M=780) & N=300\ (M=44850) & N=1000\ (M=499500)\\\hline
f(0) & 1.01\times10^{-30} & 1.40\times10^{-30} & 4.58\times10^{-30}\\
\text{parent residual} & 6.2\times10^{-13} & 9.3\times10^{-13} & 4.5\times10^{-12}\\
\text{initial growth rate (/step}) & 3.60\times10^{-2} & 2.49\times10^{-2} & 2.11\times10^{-2}\\
\text{steps per decade} & 64 & 93 & 109\\
\tau_c(f\ge0.05) & 2011 & 4844 & 7926\\
\lvert Z^{\mathsf T}Z\rvert\text{ max} & 1.4\times10^{-15} & 5.6\times10^{-15} & 2.7\times10^{-14}
\end{array}
\]

\begin{figure}
\centering
\pandocbounded{\includegraphics[keepaspectratio,alt={Large-N onset: constant-rate amplification from a seed depth of 30 decades}]{figs/dormant_growth_large_n_v1.png}}
\caption{Large-N onset: constant-rate amplification from a seed depth of
30 decades}
\end{figure}

Three readings are established.

\begin{enumerate}
\def\labelenumi{\arabic{enumi}.}
\tightlist
\item
  \textbf{The onset law holds at large \(N\).} From seeds
  \(f(0)\approx10^{-30}\), 25 decades deeper than Experiment O1
  (\(N=4\), \(\delta\ge10^{-5}\)), the dormant fraction passed straight
  through more than 20 decades of constant-factor exponential
  amplification, then left the exponential regime and entered the
  saturation region where growth slows by orders. The
  seed-depth-independent exponential onset (Section 8.2) reproduces
  unchanged at body counts 40, 300, and 1000 (about half a million
  relational waves). The implication of Axiom 0.5 --- only amplitude
  ratios act, and the arrival time of the amplification is linear in the
  seed's decade count --- is supported by direct measurement to this
  depth
\item
  \textbf{The rate depends only weakly on \(N\).} The initial growth
  rate decreases monotonically from \(3.8\times10^{-2}\) at \(N=4\)
  through \(3.6\times10^{-2}\) at \(N=40\) and \(2.5\times10^{-2}\) at
  \(N=300\) to \(2.1\times10^{-2}\) at \(N=1000\), staying within the
  same order, with the decrease slowing at large \(N\). Identifying the
  \(N\)-dependence law of the rate is a connection problem
\item
  \textbf{The measurement floor is set by the parent's construction
  accuracy.} The gentle initial rise from \(10^{-30}\) to about
  \(10^{-24}\) is a transient in which the construction residual of the
  parent eigenmode (\(\sim10^{-12}\), energy \(\sim10^{-24}\)) leaks
  into the orthogonal complement; it is distinct from the subsequent
  exponential regime. The 20+ decades above the floor are the measured
  range of dynamical amplification
\end{enumerate}

Reservation: one random seed per \(N\); no \(\delta\)-series or seed
statistics as in Experiment O1. Tracking extends to 1500 steps past
crossing (\(N=1000\): 1074 steps past crossing due to the 9000-step
cap), and the outcome classification of large \(N\) (the three classes
of Section 8.6) is outside the scope of this experiment.

\begin{center}\rule{0.5\linewidth}{0.5pt}\end{center}

\subsection{9. Discussion and
Limitations}\label{discussion-and-limitations}

\textbf{Relation to modulational instability} The onset phenomenon of
this paper shows phenomenology isomorphic to modulational instability
{[}5,6{]}: exponential growth, seed-amplitude-independent rate,
activation of discrete modes, saturation. The mechanism differs,
however. NLS-type modulational instability arises from
amplitude-dependent nonlinear coupling, whereas the generator of this
system does not depend on amplitudes (phase-only coupling, zeroth-order
homogeneous). It is not ``a kind of MI'' but ``a phase-coupled
instability with phenomenology isomorphic to MI.''

\textbf{FPUT-type recurrence} The post-saturation back-and-forth of
splitting and recombination is isomorphic to recurrence in conservative
systems {[}7,8{]}. The framework for judging recurrence versus
thermalization (equipartition \(\mathrm{PR}=M\)) follows {[}8{]}. The
bounded-mixing class shows the intermediate aspect ``recurring while the
amplitudes head toward equipartition''; refining this classification is
deferred to the sequel.

\textbf{Determinism and predictability} The collapse direction of the
\(\delta=0\) run was decided by floating-point rounding. The update rule
is fully deterministic, but the outcome depends sensitively on
microscopic differences not contained in the macroscopic description
{[}9{]}. Moreover, that the conserved global invariant \(c\) of the
initial condition separates the outcome classes at small \(N\) shows
that the value of a conserved quantity acts on outcome selection
(Section 8.7).

\textbf{Limitations} (i) The operation of increasing the component count
itself is an underdetermined freedom the axioms grant from the start
(\(X_P^2=\sum_kY_k^2\) is always solvable over the complex numbers with
a continuum of solutions) and requires no dynamical justification. That
this paper selects no particular parameterization (no explicit splitting
map) is not an omission: the physics --- conservation, freezing, outcome
selection --- appears only in the readout layer (Section 1.2).
``Splitting'' in this paper is a statement of the readout layer:
register reallocation and growth of the active-plane count within a
state space of fixed dimension. (ii) The \(\operatorname{rank}2\)
property of the parent is a numerical insight, unproven. (iii) At larger
\(N\) (some trials with \(N\ge5\)) the parent fixed-point iteration did
not converge and approximate parents were used. (iv) The classification
thresholds (\(\sigma_2>0.05\), \(\min f<0.05\), late window = second
half) are preliminary; no sensitivity analysis was performed. (v) The
origin of the saturation ceiling is unresolved. (vi) No theoretical
formula for the growth rate \(3.8\times10^{-2}\)/step (Floquet analysis
around the parent periodic orbit) has been carried out, and the
quantitative resolution of the mismatch between the initial-fit rate and
the crossing-derived effective rate (\(5.5\times10^{-2}\)) --- the
non-single-exponential nature of the growth --- also remains. (vii) The
relative activity threshold \(0.05\) is a readout convention; its
conventionality (the choice of how to count) is itself thematized in the
sequel.

\begin{center}\rule{0.5\linewidth}{0.5pt}\end{center}

\subsection{10. Connection Problems}\label{connection-problems}

\begin{enumerate}
\def\labelenumi{\arabic{enumi}.}
\tightlist
\item
  \textbf{Stopping of splitting} (sequel, second question). The
  mechanism and rate of the spontaneous \(\sigma_2\) decay in the
  \(N=4\) parent family. Identification of the destination of the decay
  (a new self-consistent configuration). Relation to external generator
  fixing (the operation that freezes all registers in the control
  experiment). The threshold of re-splitting.
\item
  \textbf{Differentiation of conserved quantities} (sequel, third
  question). Classification of the relative components generated by
  splitting via exchange symmetry.
\item
  Construction of an explicit splitting map (component-count increase,
  exact conservation of the complex square sum) and its relation to the
  growth of the active-plane count in this paper.
\item
  Floquet analysis of the parent periodic orbit and comparison with the
  measured growth rate.
\item
  The mechanism of outcome stratification by the conserved ratio \(c\)
  --- why expansion is suppressed for \(N\le4\) in the \(c=0\) stratum;
  measurement of the stratum boundary (the critical \(c\) at which
  classes switch).
\item
  Systematic measurement of the \(N\) and \(\gamma\) dependence of the
  saturation ceiling.
\item
  Proof of the parent's \(\operatorname{rank}2\) property; the existence
  problem of self-consistent parents at larger \(N\).
\item
  Statistics of selective freezing (distribution of freezing times and
  frozen allocations).
\end{enumerate}

\begin{center}\rule{0.5\linewidth}{0.5pt}\end{center}

\subsection{11. Reproducibility}\label{reproducibility}

The numerical experiments of this paper are based on the following
programs and outputs.

\subsubsection{Programs}\label{programs}

\begin{itemize}
\tightlist
\item
  \texttt{dimension\_generation\_structure/spontaneous\_splitting\_preliminary\_v1/run\_spontaneous\_splitting\_preliminary\_v1.py}
  (Experiments N, O1, scale verification, \(\delta=0\) verification)
\item
  \texttt{dimension\_generation\_structure/spontaneous\_splitting\_preliminary\_v1/run\_outcome\_classification\_preliminary\_v1.py}
  (Experiment O4)
\item
  \texttt{dimension\_generation\_structure/spontaneous\_splitting\_preliminary\_v1/run\_closure\_contrast\_v1.py}
  (Experiment O5)
\item
  \texttt{dimension\_generation\_structure/spontaneous\_splitting\_preliminary\_v1/run\_spontaneous\_splitting\_largeN\_v1.py}
  (Experiment O6)
\item
  \texttt{dimension\_generation\_structure/spontaneous\_splitting\_preliminary\_v1/run\_n\_scaling\_lowrank\_v1.py}
  (low-rank vertex-decomposition engine; dense-matrix agreement
  verification)
\item
  \texttt{dimension\_generation\_structure/spontaneous\_splitting\_preliminary\_v1/make\_largeN\_figure\_v1.py}
  (Experiment O6 figure)
\item
  \texttt{dimension\_generation\_structure/spontaneous\_splitting\_preliminary\_v1/make\_paper\_supplement\_figures\_v1.py}
  (supplementary figures)
\end{itemize}

\subsubsection{Main outputs}\label{main-outputs}

\begin{itemize}
\tightlist
\item
  \texttt{spontaneous\_splitting\_result\_v1/summary\_v1.json},
  \texttt{dormant\_fraction\_curves\_v1.csv}, 7 figures
\item
  \texttt{outcome\_classification\_result\_v1/summary\_v1.json}, 3
  figures
\item
  \texttt{closure\_contrast\_result\_v1/summary\_v1.json} (Experiment
  O5)
\item
  \texttt{largeN\_splitting\_result\_v1/fcurve\_*.csv},
  \texttt{summary\_*.json}, 1 figure (Experiment O6)
\end{itemize}

\begin{center}\rule{0.5\linewidth}{0.5pt}\end{center}

\subsection{12. Final Conclusions}\label{final-conclusions}

Dynamics preserving the closure sphere is exhausted by \(\dot X=K(X)X\)
(\(K\) skew-symmetric); the fixed generator is a special case. Merely
unfreezing the sine generator --- which was defined as a function of
phase differences from the start --- without introducing any new
coupling rule and without external forcing, the parent state decays
exponentially at an intrinsic rate independent of the seed amplitude,
while the real and imaginary parts of the complex square-sum closure are
conserved to machine precision and the allocation moves spontaneously
into dormant directions. Even the exact periodic orbit at \(\delta=0\)
collapses through the amplification of the nonzero numerical
perturbation of rounding. Just outside the boundary drawn by the no-go
of Paper 4, spontaneous splitting begins.

Existence opens without limit; only the readout saturates. The rank
bound and three-direction saturation of Paper 3, the register freezing
of Paper 4, and the outcome classification of this paper are all
statements of readout saturation, not statements of the finiteness of
existence.

\[
\boxed{
\begin{aligned}
&\text{Splitting can begin. But the outcome is not a single road.}\\
&\text{Three classes are real — expansion (establishment of a new plane, equipartition),}\\
&\text{bounded mixing (wandering of a single plane), recurrence (repeated returns toward the parent) —}\\
&\text{and the selection is decided by the }N\text{ phase structure }(N\le4\text{ no expansion, }N\ge5\text{ expansion in all trials})\\
&\text{and the conserved ratio }c.
\end{aligned}
}
\]

Scale anonymity (Axiom 0.5) was introduced as the formation rule running
through this entire construction. The unfrozen generator satisfies it
without adjustment as a zeroth-order homogeneous function; the readouts
and thresholds are all ratio quantities; and numerically all ratio
observables agreed bitwise under a \(2^{10}\)-fold scale change. The
zero square sum of Axiom 1 is the unique quadratic closure compatible
with this axiom.

The spontaneous exponential decay of \(\sigma_2\) observed in the
\(N=4\) parent family suggests that splitting can stop without an
external freezing operation. Can splitting begin? --- the answer of this
paper is yes. Can it stop? --- that is the question of the sequel.

\begin{center}\rule{0.5\linewidth}{0.5pt}\end{center}

\subsection{References}\label{references}

\subsubsection{Self-citations}\label{self-citations}

\begin{enumerate}
\def\labelenumi{\arabic{enumi}.}
\tightlist
\item
  Noriaki Kihara, ``Basic Axiom System of the Anonymous Equal-Amplitude
  Composite Wave Model v7,'' Zenodo, 2026. Version DOI:
  \href{https://doi.org/10.5281/zenodo.21469133}{10.5281/zenodo.21469133},
  Concept DOI:
  \href{https://doi.org/10.5281/zenodo.21315735}{10.5281/zenodo.21315735}.
\item
  Noriaki Kihara, ``Linear Upper Bound on Generator Rank and
  Three-Direction Saturation of Spatial-Direction Readout in N-Body
  Complete Pairwise Relational Waves,'' Zenodo, 2026. Version DOI:
  \href{https://doi.org/10.5281/zenodo.21465899}{10.5281/zenodo.21465899},
  Concept DOI:
  \href{https://doi.org/10.5281/zenodo.21465898}{10.5281/zenodo.21465898}.
\item
  Noriaki Kihara, ``Plane-Decomposition Readout in N-Body
  Fixed-Generator Systems,'' Zenodo, 2026. Version DOI:
  \href{https://doi.org/10.5281/zenodo.21468960}{10.5281/zenodo.21468960},
  Concept DOI:
  \href{https://doi.org/10.5281/zenodo.21468959}{10.5281/zenodo.21468959}.
\item
  Noriaki Kihara, ``Future-Phase Position-Acceleration Mapping and the
  Inverse-Square Law from Harmonic Closure in the Closed Two-Body AB
  Phase System v4,'' Zenodo, 2026. Version DOI:
  \href{https://doi.org/10.5281/zenodo.21468270}{10.5281/zenodo.21468270},
  Concept DOI:
  \href{https://doi.org/10.5281/zenodo.21441081}{10.5281/zenodo.21441081}.
\end{enumerate}

\subsubsection{External references}\label{external-references}

\begin{enumerate}
\def\labelenumi{\arabic{enumi}.}
\setcounter{enumi}{4}
\tightlist
\item
  T. Brooke Benjamin and J. E. Feir, ``The disintegration of wave trains
  on deep water. Part 1. Theory,'' \emph{Journal of Fluid Mechanics},
  27(3), 417--430, 1967. DOI:
  \href{https://doi.org/10.1017/S002211206700045X}{10.1017/S002211206700045X}.
\item
  Vladimir E. Zakharov and Lev A. Ostrovsky, ``Modulation instability:
  The beginning,'' \emph{Physica D}, 238(5), 540--548, 2009. DOI:
  \href{https://doi.org/10.1016/j.physd.2008.12.002}{10.1016/j.physd.2008.12.002}.
\item
  Enrico Fermi, John Pasta, and Stanislaw Ulam, ``Studies of Nonlinear
  Problems. I,'' Los Alamos Report LA-1940, 1955. DOI:
  \href{https://doi.org/10.2172/4376203}{10.2172/4376203}. (Including
  Mary Tsingou's contribution to the computational implementation.)
\item
  Gennady P. Berman and Felix M. Izrailev, ``The Fermi--Pasta--Ulam
  problem: Fifty years of progress,'' \emph{Chaos}, 15(1), 015104, 2005.
  DOI: \href{https://doi.org/10.1063/1.1855036}{10.1063/1.1855036}.
\item
  Edward N. Lorenz, ``Deterministic Nonperiodic Flow,'' \emph{Journal of
  the Atmospheric Sciences}, 20(2), 130--141, 1963. DOI:
  \href{https://doi.org/10.1175/1520-0469(1963)020\%3C0130:DNF\%3E2.0.CO;2}{10.1175/1520-0469(1963)020\textless0130:DNF\textgreater2.0.CO;2}.
\item
  Fasma Diele, Luciano Lopez, and R. Peluso, ``The Cayley Transform in
  the Numerical Solution of Unitary Differential Systems,''
  \emph{Advances in Computational Mathematics}, 8(4), 317--334, 1998.
  DOI:
  \href{https://doi.org/10.1023/A:1018908700358}{10.1023/A:1018908700358}.
\item
  David J. Thouless, ``Electrons in disordered systems and the theory of
  localization,'' \emph{Physics Reports}, 13(3), 93--142, 1974. DOI:
  \href{https://doi.org/10.1016/0370-1573(74)90029-5}{10.1016/0370-1573(74)90029-5}.
\end{enumerate}

\end{document}
